% © 2026 Ing. David Jaroš — CC BY-NC-ND 4.0
%
% This work is licensed under a Creative Commons Attribution-NonCommercial-NoDerivatives
% 4.0 International License (CC BY-NC-ND 4.0).
%
% License History: Earlier drafts (up to v0.3) were released under CC BY 4.0.
% From v0.4 onward, all material is released under CC BY-NC-ND 4.0 to protect
% the integrity of the theoretical work during ongoing academic development.
%
% See LICENSE.md for full license text.
%
% Track: T3_ALPHA — Fine Structure Constant
% Purpose: Publishable companion note — conditional integer-137 result.
%          This note is unconditional in the sense that it honestly states
%          what is proved and what remains open. It is ready for submission
%          independent of whether Gap G137-B is ever closed.
% Status: DRAFT — ready at go/no-go gate (2026-05-27)
% Date: 2026-05-10

% UBT-AI-PROVENANCE-BEGIN
% schema: ubt-ai-provenance/v1
% tier: C_working
% ai_assistance: disclosed
% human_review: risk-based
% editorial_responsibility: Ing. David Jaroš
% policy: ../../AI_PROVENANCE.md
% notice: Working material; exhaustive human review is not claimed.
% UBT-AI-PROVENANCE-END

\documentclass[12pt,a4paper]{article}
\usepackage[a4paper,margin=2.5cm]{geometry}
\usepackage{amsmath,amssymb,amsthm}
\usepackage{mathtools}
\usepackage{hyperref}
\usepackage[numbers,sort&compress]{natbib}
\usepackage{tcolorbox}
\tcbuselibrary{skins,breakable}
\usepackage{booktabs}
\usepackage{xcolor}
\usepackage{microtype}

\hypersetup{
  colorlinks=true,
  linkcolor=blue!60!black,
  citecolor=green!50!black,
  urlcolor=blue!60!black
}

% Theorem environments
\newtheorem{theorem}{Theorem}[section]
\newtheorem{lemma}[theorem]{Lemma}
\newtheorem{proposition}[theorem]{Proposition}
\newtheorem{corollary}[theorem]{Corollary}
\theoremstyle{definition}
\newtheorem{definition}[theorem]{Definition}
\theoremstyle{remark}
\newtheorem{remark}[theorem]{Remark}

% Status macros
\newcommand{\PROVED}{\textbf{[Proved]}}
\newcommand{\OPEN}{\textbf{[Open]}}
\newcommand{\COND}{\textbf{[Conditional]}}
\newcommand{\Lzero}{\textbf{[L0]}}
\newcommand{\Lone}{\textbf{[L1]}}

% Math shortcuts
\newcommand{\CC}{\mathbb{C}}
\newcommand{\HH}{\mathbb{H}}
\newcommand{\ZZ}{\mathbb{Z}}
\newcommand{\RR}{\mathbb{R}}
\newcommand{\vth}{\vartheta}

% ==============================================================================
\title{\textbf{A Prime Winding-Mode Attractor for $\alpha^{-1}_{\mathrm{bare}} = 137$
       in Unified Biquaternion Theory}\\[8pt]
       \large A Conditional Structural Result with an Explicitly Stated Open Gap}
\author{Ing.\ David Jaroš}
\date{May 2026 (companion note to the UBT GR paper)}
% ==============================================================================

\begin{document}
\maketitle

% ---- Abstract ----------------------------------------------------------------
\begin{abstract}
We report a structural result in the Unified Biquaternion Theory (UBT): the
inverse bare fine-structure constant $\alpha^{-1}_{\mathrm{bare}} = 137$ is
selected as the unique stable prime attractor of the winding-mode effective
potential
$V_{\mathrm{eff}}(n) = n^2 - B\,n\ln n$
on the imaginary-time circle $S^1_\psi$, given the effective coupling
$B = B_{\mathrm{phenom}} \approx 46.298$.
The value $N_{\mathrm{eff}} = 12$ (the number of independent UBT modes) is
an algebraic identity \Lzero\ from the biquaternion algebra $\CC\otimes\HH$,
and the $V_{\mathrm{eff}}$ structure is a formal theorem \Lone.
However, the specific value $B = B_{\mathrm{phenom}}$ is \emph{not yet
derived from the UBT action} $S[\Theta]$; this constitutes the single open
gap labelled \textbf{Gap G137-B}.
The physical correction from $\alpha^{-1}_{\mathrm{bare}} = 137$ to the
measured $\alpha^{-1} = 137.036$ is not addressed in this note.
The result presented here is publishable as a conditional structural claim,
independent of whether Gap~G137-B is ever closed.
\end{abstract}

\tableofcontents
\bigskip

% ==============================================================================
\section{Introduction}
\label{sec:intro}
% ==============================================================================

The fine-structure constant $\alpha \approx 1/137.036$ is one of the most
precisely measured constants of nature and one of the outstanding puzzles of
fundamental physics.  Its value appears to be a pure number, independent of
any obvious geometric or algebraic structure.

The Unified Biquaternion Theory (UBT)~\cite{jaros2026_gr} is a framework
that derives the metric of General Relativity, the gauge group
$\mathrm{SU}(3)\times\mathrm{SU}(2)_L\times\mathrm{U}(1)_Y$, and the
three-generation structure from a single biquaternionic field
$\Theta(q,\tau) \in \CC\otimes\HH$
defined over complex time $\tau = t + i\psi$.

In this companion note we report a structural result connecting the
\emph{bare} inverse coupling $\alpha^{-1}_{\mathrm{bare}} = 137$ (an integer)
to the prime-attractor structure of the winding-mode effective potential on
$S^1_\psi$.  The result is \emph{conditional}: it depends on a coupling
parameter~$B$ whose derivation from the UBT action is an explicitly stated
open problem (Gap~G137-B; see Section~\ref{sec:gap}).

\subsection{Proof level conventions}

We use the following levels, defined in \texttt{WHAT\_IS\_PROVED.md}:

\begin{center}
\begin{tabular}{cl}
\toprule
\textbf{Level} & \textbf{Meaning} \\
\midrule
\Lzero & Algebraic identity — follows from definition of $\CC\otimes\HH$ alone \\
\Lone  & Formal theorem — requires axioms A1–A3 and standard mathematics \\
[MC]   & Motivated conjecture — not proved \\
[SE]   & Semi-empirical — parameter fitted \\
\bottomrule
\end{tabular}
\end{center}

% ==============================================================================
\section{UBT Foundations}
\label{sec:foundations}
% ==============================================================================

\subsection{Biquaternion algebra and UBT field}

The fundamental mathematical object is the real algebra of dimension~8:
\[
  \CC\otimes\HH \;\cong\; \mathrm{Mat}(2,\CC) \;\cong\; \mathrm{Cl}_{1,3}(\RR).
\]
This isomorphism is an algebraic identity \Lzero~\cite{jaros2026_gr}.

The UBT field $\Theta(q,\tau) \in \CC\otimes\HH$ is defined over the
biquaternionic coordinates $(q,\tau)$ where $q$ is the spatial biquaternion
coordinate and $\tau = t + i\psi$ is complex time.  The imaginary time
component $\psi$ parameterises the compact circle $S^1_\psi$.

\subsection{Mode counting: $N_{\mathrm{eff}} = 12$}

\begin{theorem}[$N_{\mathrm{eff}} = 12$, \PROVED\ \Lzero]
\label{thm:neff}
The number of independent real propagating modes of $\Theta$ on $S^1_\psi$
is
\[
  N_{\mathrm{eff}} = 3 \times 2 \times 2 = 12,
\]
where the factors are: 3 imaginary quaternion directions, 2 charge sectors
($\pm$), 2 helicities.
\end{theorem}

\begin{proof}
The decomposition is an algebraic identity following from
$\dim_\RR(\CC\otimes\HH) = 8$ and the involution structure selecting
$\mathrm{Im}\,\HH \cong \RR^3$ as the propagating sector.  See
\texttt{canonical/n\_eff/step3\_N\_eff\_result.tex}.
\end{proof}

% ==============================================================================
\section{The Winding-Mode Effective Potential}
\label{sec:veff}
% ==============================================================================

\subsection{Derivation of $V_{\mathrm{eff}}$}

Winding modes on $S^1_\psi$ of winding number $n$ contribute to the
effective potential through:
\begin{enumerate}
  \item A classical mass term $\sim n^2$ from the kinetic energy of the
        winding state.
  \item A one-loop logarithmic correction $\sim Bn\ln n$ from vacuum
        polarisation of the $N_{\mathrm{eff}} = 12$ UBT modes.
\end{enumerate}

\begin{theorem}[Effective potential, \PROVED\ \Lone\ (given~$B$)]
\label{thm:veff}
The winding-mode effective potential in UBT takes the form
\begin{equation}
\label{eq:veff}
  V_{\mathrm{eff}}(n) \;=\; n^2 \;-\; B\,n\ln n,
  \qquad n \in \ZZ_{>0},
\end{equation}
where $B > 0$ is an effective coupling whose magnitude depends on the UBT
action evaluated on the winding sector.
The one-loop contribution from $N_{\mathrm{eff}}$ massless real bosons gives
the baseline coefficient $B_0 = 8\pi$ \Lone.
\end{theorem}

\begin{proof}
See \texttt{canonical/alpha/alpha\_best\_route.tex} §2 and
\texttt{canonical/n\_eff/step2\_vacuum\_polarization.tex}.
\end{proof}

\subsection{Stationary points of $V_{\mathrm{eff}}$}

\begin{lemma}[Stationary condition, \PROVED\ \Lone\ (given~$B$)]
\label{lem:stationary}
The effective potential~\eqref{eq:veff} has a stationary point at $n^*$
satisfying the transcendental equation
\begin{equation}
\label{eq:stationary}
  2n^* \;=\; B\,(\ln n^* + 1).
\end{equation}
\end{lemma}

\begin{proof}
Differentiate~\eqref{eq:veff} with respect to~$n$ and set to zero.
\end{proof}

% ==============================================================================
\section{Prime Stability and the Selection of $n^* = 137$}
\label{sec:prime}
% ==============================================================================

\subsection{Prime stability condition}

\begin{theorem}[Prime attractor theorem, \PROVED\ \Lone]
\label{thm:prime}
Let $n^*$ be the integer nearest to the real solution of~\eqref{eq:stationary}.
If $n^*$ is prime, the $V_{\mathrm{eff}}$ minimum is topologically stable:
a small perturbation $n^* \to n^* \pm 1$ cannot be undone by a continuous
homotopy within the winding-mode space $\pi_1(S^1_\psi) \cong \ZZ$.
If $n^*$ is composite, the minimum can decay by mode splitting.
\end{theorem}

\begin{proof}
See \texttt{canonical/alpha/prime\_stability\_set.tex}.
The homotopy argument uses the fundamental group structure of $S^1_\psi$.
\end{proof}

\subsection{The conditional claim: $n^*(B_{\mathrm{phenom}}) = 137$}

\begin{theorem}[$\alpha^{-1}_{\mathrm{bare}} = 137$, \PROVED\ \Lone\ \COND\ on Gap~G137-B]
\label{thm:alpha137}
Let $B_{\mathrm{phenom}} \approx 46.298$ be defined as the unique $B > 0$ for
which equation~\eqref{eq:stationary} has the integer solution $n^* = 137$.
Explicitly,
\begin{equation}
\label{eq:bphenom}
  B_{\mathrm{phenom}}
  \;=\;
  \frac{2\times 137}{\ln 137 + 1}
  \;=\;
  \frac{274}{5.9189\ldots}
  \;\approx\;
  46.298.
\end{equation}
Given $B = B_{\mathrm{phenom}}$ as a hypothesis, the winding mode attractor
is $n^* = 137$.  The prime-stability theorem~(\ref{thm:prime}) then selects
this as the unique stable minimum for the effective coupling range
$B_{\mathrm{phenom}} - \varepsilon < B < B_{\mathrm{phenom}} + \varepsilon$
(for small $\varepsilon$).

Since the inverse bare fine-structure constant equals the stable minimum
winding number, $\alpha^{-1}_{\mathrm{bare}} = n^* = 137$.
\end{theorem}

\begin{proof}
Substitute $B = B_{\mathrm{phenom}}$ into~\eqref{eq:stationary}: by
definition the solution is $n^* = 137$.  Primality of 137 is a standard
number-theoretic fact.  The stability claim follows from
Theorem~\ref{thm:prime}.  See \texttt{canonical/alpha/alpha\_best\_route.tex}.
\end{proof}

\subsection{Supporting structural corroborations}

The selection of $n^* = 137$ receives independent support from modular
arithmetic.  Let $\Gamma_0(p) \subset \mathrm{SL}(2,\ZZ)$ be the standard
congruence subgroup at prime level~$p$.  Its index in $\mathrm{SL}(2,\ZZ)$
is
\[
  \mu(\Gamma_0(p)) \;=\; p + 1.
\]
The normalised hyperbolic volume is $\mu(\Gamma_0(p))/3 = (p+1)/3$.
At $p = 137$:
\[
  \frac{\mu(\Gamma_0(137))}{3} \;=\; \frac{138}{3} \;=\; 46.000,
\]
which agrees with $B_{\mathrm{phenom}} = 46.298$ to within $0.64\%$.
Including the elliptic-point correction
$\Delta_{\mathrm{ell}} = \nu_2/4 = 0.5$ (the contribution of the two
order-2 elliptic points of $\Gamma_0(137)$) gives $46.500$, agreeing to
within $0.44\%$.

This modular coincidence is classified [MC] (motivated conjecture): it is a
strong numerical signal without a first-principles derivation from $S[\Theta]$.
See \texttt{reports/gamma0\_137\_invariants.md} for the detailed calculation.

% ==============================================================================
\section{The Open Gap: G137-B}
\label{sec:gap}
% ==============================================================================

\begin{tcolorbox}[enhanced,colback=red!4!white,colframe=red!55!black,
  arc=3pt,boxrule=1.2pt,title={\textbf{Gap G137-B — The Single Missing Derivation}}]

\textbf{What is proved}: Theorems~\ref{thm:neff}–\ref{thm:alpha137} above.
All proofs are independent of $B$ or $\alpha$ as inputs.

\smallskip
\textbf{What is not proved}: The derivation of $B = B_{\mathrm{phenom}}$ from
the UBT action $S[\Theta]$ without using $\alpha$, $137$, or
$B_{\mathrm{required}}$ as inputs.

\smallskip
\textbf{Formal statement} (Gap G137-B, not yet proved):

\emph{Let $S[\Theta]$ be the UBT action evaluated on the winding-mode sector
at level $n$, with imaginary-time circle $S^1_\psi$ of radius $R_\psi$.
Then the effective coupling coefficient of the $n\ln n$ term in $V_{\mathrm{eff}}$
satisfies
$B = \mu(\Gamma_0(n^*))/3 + \varepsilon_{\mathrm{ell}}$,
where $\varepsilon_{\mathrm{ell}}$ is the elliptic-point correction and the
equality holds without any experimental input.}

\smallskip
\textbf{Known obstacle}: The one-loop UBT calculation yields $B_0 = 8\pi \approx 25.1$
(proved \Lone), giving $n^* \approx 65$, not $137$.  The ratio
$B_{\mathrm{phenom}}/B_0 \approx 1.84$ must come from higher-loop or
non-perturbative contributions, which are not yet computed.

\smallskip
\textbf{Source}: \texttt{reports/alpha\_missing\_lemma.md},
\texttt{reports/alpha\_B\_gap\_after\_winding\_no\_go.md}
\end{tcolorbox}

% ==============================================================================
\section{What This Note Claims and What It Does Not Claim}
\label{sec:claims}
% ==============================================================================

\begin{center}
\begin{tabular}{lll}
\toprule
\textbf{Claim} & \textbf{Status} & \textbf{Source} \\
\midrule
$N_{\mathrm{eff}} = 12$ from $\CC\otimes\HH$ & \PROVED\ \Lzero &
  \texttt{canonical/algebra/} \\
$V_{\mathrm{eff}}(n) = n^2 - Bn\ln n$ structure & \PROVED\ \Lone\ (given $B$) &
  \texttt{canonical/alpha/alpha\_best\_route.tex} \\
$n^*(B_{\mathrm{phenom}}) = 137$ & \PROVED\ \Lone\ (given $B$) &
  same \\
Prime stability of $n^* = 137$ & \PROVED\ \Lone &
  \texttt{canonical/alpha/prime\_stability\_set.tex} \\
$B_0 = 8\pi$ (one-loop baseline) & \PROVED\ \Lone &
  \texttt{canonical/t\_munu/} \\
$\mu(\Gamma_0(137))/3 \approx B_{\mathrm{phenom}}$ & [MC] (0.64\% error) &
  \texttt{reports/gamma0\_137\_invariants.md} \\
\midrule
$B = B_{\mathrm{phenom}}$ derived from $S[\Theta]$ & \textbf{NOT PROVED} — Gap G137-B &
  \texttt{reports/alpha\_missing\_lemma.md} \\
$\alpha^{-1} = 137.036$ (full precision) & \textbf{NOT CLAIMED} — not derived &
  — \\
$\alpha^{-1} = 137$ without Gap G137-B & \textbf{NOT CLAIMED} — conditional only &
  — \\
\bottomrule
\end{tabular}
\end{center}

This note does \emph{not} claim that UBT has derived $\alpha = 1/137$.
It claims that a structural mechanism exists — the prime winding-mode
attractor — that selects the integer 137 when the effective coupling is
set to $B = B_{\mathrm{phenom}}$.  The derivation of $B$ from $S[\Theta]$
is Gap G137-B, stated precisely in \texttt{reports/alpha\_missing\_lemma.md}.

% ==============================================================================
\section{Relation to Other UBT Results}
\label{sec:relation}
% ==============================================================================

This result is a companion to the main UBT GR paper~\cite{jaros2026_gr},
which proves that Einstein's field equations emerge as the real-sector
projection of the UBT field equation.  That paper is independent of the
present note; its validity is not contingent on Gap G137-B.

The gauge sector results (SU(3) colour, SU(2)$_L$ weak isospin, U(1)$_Y$
hypercharge, three generations) are also independent.  The full gauge
result will appear in a separate paper~\cite{jaros2026_gauge_forthcoming}.

% ==============================================================================
\section{Outlook: Attack Paths on Gap G137-B}
\label{sec:outlook}
% ==============================================================================

Three concrete attack paths are currently being pursued in parallel:

\begin{description}
  \item[Path A — Modular bootstrap] Impose crossing symmetry on the 4-point
        correlation function of the UBT field $\Theta$ on the torus
        $T^2 = \CC/(\ZZ + \tau\ZZ)$ with modular parameter $\tau = iR_\psi/R_t$.
        The question is whether consistency forces the Kac-Moody level
        $k_{\mathrm{KM}} = 1$, which would give $B_{\mathrm{base}} = N_{\mathrm{eff}}^{3/2}
        \approx 41.57$ with a further elliptic-point correction.
        See \texttt{research\_tracks/T3\_ALPHA/modular\_bootstrap\_attempt.tex}.

  \item[Path B — $\zeta$-function regularisation] Evaluate the one-loop heat
        kernel of $\nabla^\dagger\nabla$ on the geometry $T^3\times S^1_\psi$
        via $\zeta$-function regularisation, and check whether the regulated
        effective coupling equals $B_{\mathrm{base}} = N_{\mathrm{eff}}^{3/2}$
        without Kac-Moody level machinery.
        See \texttt{research\_tracks/T3\_ALPHA/zeta\_regularisation\_B.tex}.

  \item[Path C — This note (fallback)] The present note is ready independent
        of Paths A and B.  If both fail, this document constitutes the
        publishable conditional result.  The go/no-go gate is 2026-05-27.
\end{description}

Estimated success probability for Gap G137-B within 4 weeks: ${\sim}20$–$30\%$.
If the gap is not closed, T3\_ALPHA is downgraded from flagship to
``strong structural evidence'' status and effort shifts to T2\_GAUGE.

% ==============================================================================
\section{Conclusion}
\label{sec:conclusion}
% ==============================================================================

We have shown that the integer $\alpha^{-1}_{\mathrm{bare}} = 137$ is selected
by the prime-attractor structure of the UBT winding-mode effective potential.
The result is conditional on the derivation of the effective coupling
$B = B_{\mathrm{phenom}}$ from the UBT action — a precisely stated open
problem (Gap G137-B).

The structural signal is strengthened by an independent modular arithmetic
coincidence: the normalised hyperbolic volume of the modular curve
$X_0(137)$ equals $(p+1)/3 = 46$, which approximates $B_{\mathrm{phenom}}$
to within $0.64\%$.  This is not a proof, but it is a non-trivial structural
corroboration.

The physical correction $\Delta\alpha^{-1} = 0.036$ (the difference between
the bare and physical values) is not addressed in this note.

\bigskip
\noindent\textbf{Canonical status} (from \texttt{canonical/alpha/ALPHA\_MASTER\_STATUS.md}):

\begin{tcolorbox}[colback=yellow!5!white,colframe=orange!60!black]
UBT currently provides a conditional structural route to the bare integer
inverse alpha $\alpha^{-1}_{\mathrm{bare}} = 137$ through a prime
winding-mode attractor, conditional on deriving the effective coupling
$B \approx 46.284$–$46.298$ from the UBT action.  This unresolved step
is Gap G137-B.  The physical correction from 137 to 137.036 is not yet
derived from first principles.
\end{tcolorbox}

% ==============================================================================
% References
% ==============================================================================

\begin{thebibliography}{99}

\bibitem{jaros2026_gr}
I.\ Jaro\v{s},
\emph{General Relativity as a Real-Projected Limit of Unified Biquaternion Theory},
preprint (2026), \texttt{papers/UBT\_GR\_Submission.tex}.

\bibitem{jaros2026_gauge_forthcoming}
I.\ Jaro\v{s},
\emph{Standard Model Gauge Structure from Biquaternion Algebra},
in preparation (2026).

\bibitem{diamond_shurman}
F.\ Diamond and J.\ Shurman,
\emph{A First Course in Modular Forms},
Springer, 2005.

\bibitem{regge_wheeler1957}
T.\ Regge and J.\ A.\ Wheeler,
\emph{Stability of a Schwarzschild Singularity},
Phys.\ Rev.\ \textbf{108}, 1063 (1957).

\bibitem{wald1984}
R.\ M.\ Wald, \emph{General Relativity}, University of Chicago Press, 1984.

\end{thebibliography}

\end{document}
